% Marca de capítulo con hipervínculo
\makeatletter
\def\chaptermark#1{\markright{\protect\hyper@linkstart{link}{\@currentHref}{\thechapter.~#1}\protect\hyper@linkend}}
\makeatother
\renewcommand{\sectionmark}[1]{}

% Estructura encabezado
\setlength{\headheight}{30pt}
% Compensa los 14 pt añadidos al encabezado para conservar la posición del texto.
\addtolength{\topmargin}{-14pt}
\newlength{\encabezadocapituloancho}
\settowidth{\encabezadocapituloancho}{\scshape\abvMateria}
\setlength{\encabezadocapituloancho}{\dimexpr.5\textwidth-.5\encabezadocapituloancho-1em\relax}
\renewcommand{\headrulewidth}{0.6pt} % Línea de encabezado
\pagestyle{fancy}
\fancyhf{} % Limpia marcas de encabezado y pie
\fancyhead[l]{\fontfamily{lmr}\selectfont\abvAlumno}
\fancyhead[r]{\fontfamily{lmr}\selectfont\parbox[b]{\encabezadocapituloancho}{\raggedleft\rightmark}}
\fancyhead[c]{\scshape\abvMateria}
\fancyfoot[c]{\hyperlink{toc}{\thepage}}

% Aplica estilo vacío solo a la página en blanco
\let\origcleardoublepage\cleardoublepage
\renewcommand{\cleardoublepage}{%
    \clearpage
    {\pagestyle{empty}\origcleardoublepage}
}
